\chapter{Source and Evidence Map}
\label{app:evidence}
The source examined is commit \code{f7d0aed594a1589ff1fe24c2cb0222c149abe322} of the project repository \citep{pytutorrepo}. We checked it against the remote commit during outline preparation. Paths below refer to that repository unless marked as thesis-local. Result files from earlier runs are identified separately; a fresh full-system evaluation was not performed for every artefact.

\begin{longtable}{p{46mm}p{103mm}}
\caption{Implementation sources used for the manuscript.}\label{tab:sources}\\
\toprule
Mechanism & Source path\\
\midrule
\endfirsthead
\toprule Mechanism & Source path\\\midrule
\endhead
Diagnosis routing & \path{backend/app/services/misconception_service.py}\\
Structural orchestration & \path{backend/app/analysis/ast_analyzer.py}\\
Category visitors & \path{backend/app/analysis/ast_visitors.py}\\
Evidence scoring & \path{backend/app/analysis/confidence.py}\\
Observable error classes & \path{backend/app/analysis/error_class.py}\\
Category catalogue & \path{backend/app/data/misconceptions.json}\\
Teaching policy & \path{backend/app/services/pedagogy_service.py}\\
Chat integration & \path{backend/app/services/chat_service.py}\\
Exercise integration & \path{backend/app/api/v1/exercises.py}\\
Profile and recommendations & \path{backend/app/services/profile_service.py}\\
Optional time decay & \path{backend/app/services/profile_decay_service.py}\\
Retrieval integration & \path{backend/app/services/rag_service.py}\\
Execution boundary & \path{backend/app/sandbox/runner.py}\\
Feature defaults & \path{backend/app/core/config.py}\\
\bottomrule
\end{longtable}

\section{Historical Evaluation Artefacts}
The diagnosis files are \path{evaluation/v2_eval_results.json} and \path{evaluation/v2_holdout_results.json}. The four input files are \path{v2_test_cases.json}, \path{clean_code_cases.json}, \path{v2_holdout_cases.json}, and \path{holdout_clean_cases.json} in the same directory. The historical with-question prediction for every positive case was reproduced by the thesis component recheck.

The retrieval files are \path{evaluation/rag_eval_results_tfidf.json}, \path{evaluation/rag_eval_results_vector.json}, and \path{evaluation/rag_eval_results_hybrid.json}. Query labels are in \path{evaluation/rag_golden.json}. The exploratory explanation experiment is documented in \path{research/faithfulness/RESULTS.md} and the two \path{results_real_deepseek_v2_run1.json} and \path{results_real_deepseek_v2_run2.json} files. The reported observations use these real-model outputs, excluding the stub-model demonstration.

\section{Recheck Procedure}
The thesis-local \path{evidence/recheck_components.py} script takes an optional \code{--repo} path. It loads the original visitors and analyser, compiles the original text-match function, supplies the existing input cases, and saves per-case predictions in \path{evidence/component_recheck.json}. The JSON also records the Python version, source commit, input hashes, and class counts. The structural with-question predictions are asserted equal to the historical files. The M3 example and its evidence score are checked separately.

\begin{lstlisting}
python evidence/recheck_components.py --repo PATH_TO_REPOSITORY
\end{lstlisting}
The command reproduces component predictions. It leaves external-model runs, database state, application tests, and historical timing outside its scope. The logging shim runs only within the recheck process and changes no production source file.

\chapter{Additional Evaluation Details}
\section{Category Support}
\begin{table}[htbp]
\centering
\caption{Number of positive cases by category. Clean controls are counted separately.}
\begin{tabular}{lrrrrrrrr}
\toprule
Set & M1 & M2 & M3 & M4 & M5 & M6 & M7 & M8\\
\midrule
Development & 4 & 3 & 7 & 7 & 5 & 6 & 8 & 7\\
Holdout & 4 & 4 & 4 & 4 & 4 & 4 & 3 & 3\\
\bottomrule
\end{tabular}
\end{table}

\section{Boundary Examples}
The M3 visitor identifies the direct assignment in the first example. The other forms lack that node pattern. These examples examine the visitor alone; a service with further fallback handling may return a different result.
\begin{lstlisting}[language=Python]
# Direct assignment from an attribute call: M3 pattern.
new = items.append(5)

# Standalone mutation: not the assignment pattern.
items.append(5)

# Annotated assignment: AnnAssign, not Assign.
new: object = items.append(5)

# Wrapped call: assignment value is a List, not Call.
new = [items.append(5)]
\end{lstlisting}
Even a direct assignment can have a different meaning for a custom object whose \code{append} method returns a value. Receiver type and task purpose are therefore separate questions from the node relationship.

\section{Exploratory Automatic Review}
The historical multi-judge report covers ten examples, with a mean score of 2.40, revision rate of 9/10, and mean inter-judge agreement of approximately 0.405. There is no human reference. These values describe the automated review process: agreement between judges leaves agreement with educators untested, and the revision count does not show whether revisions improved a reply. The pilot is insufficient to validate teaching quality.

\chapter{Reproducibility and Interpretation Checklist}
The following distinctions govern reuse of the results:
\begin{enumerate}
\item Development cases validate current coverage; they do not provide independent generalisation evidence.
\item Holdout cases were separately authored within the project, not collected from an external learner population.
\item Structural diagnosis with a benchmark question has different information from the native text baseline and from current chat routing.
\item The \code{ast-only} component includes targeted handling of parse errors and excludes the LLM service fallback.
\item A primary label is selected by rule priority before evidence-score conversion; it is not a validated cognitive root cause.
\item Code defaults, documented demonstration settings, stored experimental settings, and a deployed environment are distinct configurations.
\item Retrieval scores do not measure final-answer faithfulness; automatic review does not measure learning gains.
\item The current preparation rechecked components, not the full historical backend suite or a fresh external-model run.
\end{enumerate}
The accompanying editing notes list the evidence still needed before final submission. They contain no additional experimental results.
